\chapter{Synthesizing Findings into a Pilot Report (Day 12)}
\label{chap:day12}

\section{Motivation}

Chapters~\ref{chap:day01}--\ref{chap:day11} each ran, or rolled up, one stage of the pilot ---
environment setup, sample selection, baseline inference, region extraction, and the six metrics
in turn --- and reported its own results against its own evidence. None of those chapters had
the job of stepping back and asking what the eleven days look like \emph{as a single body of
work}, written for a reader (Professor Mustafa Abdallah) who was not present for any individual
day and needs one document rather than eleven. Day 12's task was exactly that: produce a
4--6 page technical memo synthesizing Days 1--11 for the upcoming collaboration meeting ---
motivation, dataset, model, methods, the six metrics' headline results, cross-cutting findings,
limitations, and next steps --- ahead of, and separate from, this chapter-by-chapter report
itself. The artifact produced is \texttt{pilot/report/pilot\_report\_draft.md} (commit
\texttt{c2a1f85}).

\section{The Report's Structure and Editorial Policy}

The memo's structure follows the natural order of the pilot itself rather than a generic report
template: Motivation $\rightarrow$ Dataset $\rightarrow$ Model and the grounding-proxy adaptation
$\rightarrow$ Methods summary (one table mapping each of the six metrics to its method)
$\rightarrow$ Results (one headline table) $\rightarrow$ Cross-cutting findings $\rightarrow$
Limitations $\rightarrow$ Next steps and the specific points where Professor Abdallah's guidance
is wanted. This is the same structure this chapter set later expanded, day by day, into
Chapters~\ref{chap:day01}--\ref{chap:day11}.

One editorial policy adopted for the memo is worth reporting explicitly as a methodological
choice in its own right, not just as a stylistic habit. The memo's own header states it
verbatim: \textbf{``Every number below traces to a specific CSV or figure produced in Days 1-11
(\texttt{pilot/results/}); none are invented or estimated for this writeup.''} Every figure
quoted in the memo --- the 163-sample actual n, the per-rule sensitivity numbers, the six
metrics' headline and per-class values, the $\sim$0.5 GPU-hour efficiency estimate --- was
checked against the specific results file that produced it before being written down, rather
than recalled from memory or rounded for readability. This discipline costs time under a
14-day timebox, where a faster memo could have been drafted from the day-by-day findings
narratives alone; it was kept anyway, because a synthesis document that mixes verified and
remembered numbers without flagging which is which is exactly the kind of artifact that erodes
trust once a single number turns out to be wrong.

\section{Cross-Cutting Synthesis as Its Own Activity}

The memo's most distinct contribution is not any single new measurement --- Day 12 ran no new
model calls and computed no new metric --- but the act of reading the eleven days side by side
and noticing where the \emph{same} mechanism surfaces independently in more than one of them.
\textbf{Noticing a mechanism recur across independently-run days is a different activity from,
and adds information beyond, any single day's own finding}: a day-by-day report tells the reader
what happened on each day; a synthesis tells the reader which of those eleven stories are
actually one story told from several angles. This is the activity that turns eleven days of
isolated results into a coherent picture, and it is why Day 12 exists as a separate step rather
than being folded into Day 11's already-completed rollup.

Two recurring mechanisms, each first identified independently within a single day's own
analysis and only connected explicitly during this synthesis pass, illustrate the activity:

\textbf{The fallback-rerouting mechanism (Chapters~\ref{chap:day05}, \ref{chap:day09},
\ref{chap:day10}).} Masking the worker region does not only remove visual evidence; for
samples where the worker box was the only detected box, it also removes the precondition the
person-relative check needs to run at all, silently rerouting the proxy to a cruder scene-level
fallback whose answer depends on a different signal entirely. Chapter~\ref{chap:day05} first
surfaced this as the explanation for an otherwise paradoxical result (masking \emph{more}
regions producing a slightly \emph{lower} flip rate). Chapter~\ref{chap:day09} found the same
mechanism again, independently, as a driver of answer-flips under occlusion. Chapter~\ref{chap:day10}
found it a third time, independently, inflating a small slice of completeness verdicts. No
single one of those three chapters claims to have discovered a general property of the
pipeline; together, they show one.

\textbf{The area-based region-ranking confound (Chapters~\ref{chap:day04}, \ref{chap:day06},
\ref{chap:day07}, \ref{chap:day10}).} Chapter~\ref{chap:day04} ranks candidate regions by raw
pixel area, by design, with no detection-class signal available to do otherwise. That single
design choice then resurfaces as the explanation for the weakest result in three separate,
independently-run metric chapters: Chapter~\ref{chap:day06}'s sparsity measurements concentrate
mechanically for small objects regardless of reasoning quality; Chapter~\ref{chap:day07}'s
stability IoU is markedly lower for the small hard-hat box than for the large excavator box;
and Chapter~\ref{chap:day10}'s completeness rate is lower for rule\_1 (PPE, small object) than
for rule\_3 (guardrail, often the largest object in frame). Each of the three metric chapters
reported its own number as that metric's own finding; only reading all four chapters together
makes clear that one upstream ranking heuristic, not three unrelated weaknesses, is responsible.

\section{Standard vs. Non-Standard Choices}

\textbf{Standard}: producing a structured technical report that synthesizes a multi-stage
empirical pilot --- motivation, methods, results, limitations, next steps --- ahead of a
collaboration meeting is itself an entirely ordinary practice, and the memo's section ordering
follows that convention closely.

\textbf{Non-standard, and worth naming explicitly}: the strict ``every number must trace to a
specific file'' sourcing discipline described above, maintained even under a 14-day timebox
where it would have been easy, and faster, to write the memo from recollection of each day's
findings narrative rather than re-opening each day's CSV or figure to confirm the exact value
before quoting it. Most technical memos written under this kind of deadline pressure relax that
discipline somewhere; this one did not.

\section{Implications for the Paper}

The process by which the cross-cutting findings above were produced matters as much as the
findings themselves, for the paper's methodology section.
Findings were tracked day by day, inside each day's own chapter, at the time each day's work was
done --- the fallback-rerouting mechanism was written down as a Day 5 finding before Day 9 or
Day 10 existed, and the area-ranking confound was written down as a Day 4 design note before
Day 6, Day 7, or Day 10 existed. The cross-cutting connections were drawn only afterward, during
this Day 12 synthesis pass, by rereading the already-written day-by-day findings side by side.
No day's findings narrative was revised after the fact to make a later day's discovery look like
it had been anticipated. This ordering is itself a methodological safeguard worth stating in the
paper: a synthesis that retrofits early findings to match a later, more complete picture risks
manufacturing the appearance of a coherent story where the evidence trail does not actually
support one. Reporting the two recurring mechanisms above as genuine cross-chapter discoveries,
rather than as if they had been the plan from Day 1, is only credible because the day-by-day
record this report is built from shows exactly when each piece was actually found.

This also sets the template the paper should follow: a day-by-day (or stage-by-stage) results
narrative, written contemporaneously, followed by a separate synthesis pass that explicitly
distinguishes single-stage findings from cross-stage patterns, rather than a single flat results
section that blends the two. The pilot's own evidence for why that separation matters is the two
mechanisms above --- neither would have been visible from any one stage's results table alone.
